\section{Relation to Previous Research}


The results of this thesis refine and support several findings discussed in
Chapter~\ref{ch:rel_work}, but they also show that the relationship between
similarity-based and execution-based evaluation is more complex than a simple
increase in BLEU score would suggest.

First, the observed effect of fine-tuning is consistent with previous work on
task-specific adaptation of code models. Liu et al. showed that
parameter-efficient fine-tuning can improve model performance in code change
learning tasks \cite{peft_code}. Although their work focuses on code change
learning rather than Java-to-C++ translation specifically, the tendency observed
in this thesis is similar: adapting a pre-trained model to a specific code-related
task improves its performance. In the present experiments, fine-tuning increased
the average BLEU score of \texttt{starcoder2-3b} from 0.122 to 0.189 and
increased the number of accepted solutions from 55 to 133. At the same time,
compilation errors decreased from 390 to 287. The same pattern was observed for
\texttt{Qwen2.5-Coder-3B-Instruct}: the average BLEU score increased from
0.178 to 0.217, accepted solutions increased from 101 to 168, and compilation
errors decreased from 324 to 219. These results therefore agree with the general
claim that task-specific fine-tuning improves code model performance, and they
extend this observation to the Java-to-C++ translation setting used in this thesis.

Second, the stronger performance of the Qwen model is consistent with the
model-level results reported for Qwen2.5-Coder. Hui et al. describe
Qwen2.5-Coder as a code-specific model trained on a large programming corpus
and evaluated on standard programming benchmarks such as HumanEval and MBPP
\cite{qwen_coder}. In the present study, Qwen also outperformed StarCoder2 in
the Java-to-C++ translation task. The base Qwen model achieved a higher average
BLEU score than the base StarCoder2 model, and it also produced more accepted
solutions. The same ordering remained after fine-tuning, where the fine-tuned Qwen
model achieved the highest average BLEU score and the highest number of accepted
solutions among all evaluated models. This does not mean that the present
experiment reproduces the HumanEval or MBPP benchmark results directly, because
the task and evaluation setup are different. However, it does show that the stronger
general coding performance reported for Qwen2.5-Coder also appears in this
translation-specific evaluation.

Third, the correlation analysis performed in this thesis further supports earlier
criticism of BLEU in code-related evaluation. Previous research has shown that
BLEU does not always reflect the semantic correctness of generated or translated
code, because functionally equivalent programs may have different token sequences,
while lexically similar programs may still contain logical errors
\cite{tran2019ruby, evtikhiev2022bleu}. The results of this thesis are consistent
with this view. BLEU is useful for capturing broad similarity trends, and the
statistically significant correlations show that it is not meaningless. Nevertheless,
the moderate correlation values demonstrate that BLEU alone is insufficient for
evaluating code translation quality. Therefore, the present findings agree with
previous studies in treating BLEU as an auxiliary similarity metric rather than as
a standalone measure of functional correctness.

Finally, the execution-based results agree with recent work that emphasizes the
importance of compiling and running generated code. CodeEditorBench was
introduced to evaluate code editing capabilities of large language models, including
translation, in practical programming scenarios \cite{codeeditorbench}. The results
of this thesis confirm the value of such an evaluation setup. Even when BLEU
improved after fine-tuning, execution-based evaluation still revealed many failures
that would not be visible from textual similarity alone. The best-performing model,
the fine-tuned Qwen variant, produced 168 accepted solutions, but it still produced
219 compilation errors and 96 wrong answers. This confirms the gap identified in
earlier research between surface-level similarity and actual program correctness.
In this respect, the present study agrees with the broader direction of recent code
evaluation research: similarity-based metrics can provide useful auxiliary
information, but functional correctness in code translation requires execution-based
evaluation whenever test cases are available.